% ===== §2 Symptomatology: The Visible Forms of Damage =====
\section{Symptomatology: How Power Damages Generativity}
\label{sec:symptom}

\noindent
Before a thing can be explained it must be described. This section sets out the \emph{symptoms}---the observable forms in which power damages the generativity of an intimate bond---and holds itself, as far as it can, to description. The mechanisms that produce these forms are the matter of the etiology that follows (\xref{sec:encoding}--\xref{sec:ir}); here the task is only to say what is seen, and to draw the few distinctions that will let the later diagnosis proceed without confusion. A symptomatology is not yet a pathology. Much of what it records will turn out, on examination, not to be disease at all.

\subsection{Three presentations}
\label{sub:symptom-three}

Power announces its damage to a bond in three characteristic ways. They are not three diseases but three faces, and a single underlying disturbance may show now one, now another.

\subsubsection{Estrangement (\emph{surechigai})}

The first and most familiar is the one the Prelude named: the passing-without-meeting. Two subjects who remain bound, who have not quarrelled and have not betrayed, find nonetheless that they no longer arrive at one another. A gesture of care is offered and is not received as care; a silence that one reads as peace the other reads as withdrawal; the rhythms of two lives, once entrained, drift out of phase. Phenomenologically, estrangement is marked by what it is \emph{not}: it is not collision but near-miss; not the heat of conflict but the cool of trajectories that no longer cross. In the geometric language the series has developed elsewhere, it is the condition of \emph{zero-phase adjacency}---two paths brought close in space while the phase that would have accumulated between them goes to nothing.\footnote{For the holonomic apparatus---the positive geometric phase of a good cycle, and its arrest---see \xref{sec:phase} below, and \citet{paperix}. The present section uses the vocabulary descriptively; its formal content is deferred.} What is felt is distance without departure: still together, no longer meeting.

\subsubsection{Domination}

The second presentation is the one the word \emph{power} most readily calls to mind, and it is the form to which the bulk of this paper's argument is addressed. Here the disturbance is not that the two pass one another but that one of them is, increasingly, absorbed into the other's path. The signs are well known and easy to enumerate: decisions that were once joint become unilateral; one subject's preferences come to set the default against which the other's must justify themselves; the range over which one party may act, dissent, or simply be without consequence narrows. Domination need not be cruel and is rarely announced. Its mildest and most common form is the quiet contraction of one subject's world---fewer friends of one's own, fewer projects of one's own, fewer judgments ventured without first consulting the other's likely response---until the relation that was a coupling of two centres has become the orbit of one body about another. The symptom, stated plainly, is the conversion of a difference between two subjects into a standing asymmetry that no longer moves.

\subsubsection{The decoherence of trust}

The third presentation is the subtlest, and it furnishes the bridge to the theory of trust the series has already built. A bond may retain every outward form of trust---the vows still spoken, the routines still kept, the signals still exchanged---while the thing those forms were meant to carry has drained away. This is not the \emph{absence} of trust, which would be plain; it is trust \emph{hollowed}, the symbolic vessel intact and the real content gone. The series has a name for the limiting case: \emph{forged trust}, in which the apparatus of fidelity is perfectly maintained over a void.\footnote{See \citet{paperxiii}, on forged trust and the three channels of trust-formation; and \xref{sec:phase}--\xref{sec:ir} below for the mechanisms by which the real content drains while the symbolic shell persists.} Decoherence is the process by which a living trust is brought toward that limit---each act of fidelity performed a little more for the form of it, a little less from the coupling it once expressed---and its phenomenology is the uncanny one of a relation in which nothing can be pointed to as wrong and nothing is any longer right.

\subsection{The cardinal distinction: reparable misalignment versus structural shadow}
\label{sub:symptom-distinction}

Across all three presentations runs a single distinction on which the entire later argument will turn, and which the symptomatology must establish before any mechanism is named. Not every symptom is a pathology.

\begin{claim}[Reparable versus structural]
A symptom of damage may be \emph{reparable misalignment}---a local, contingent failure of meeting that the bond can metabolise and from which it can return---or it may be \emph{structural shadow}: a form of estrangement, asymmetry, or hollowing that belongs to the very constitution of a generative bond and cannot be removed without removing the generativity itself. The first calls for repair. The second calls for recognition. To treat the second as though it were the first---to attack, as a fault to be corrected, what is in fact a necessary moment of the bond's own life---is itself a principal cause of damage.
\end{claim}

\noindent
This distinction is not a hedge. It is the hinge of the paper's first claim (\xref{sub:intro-claims}): that estrangement is, in the first instance, not a malfunction but a necessary moment of self-maintenance. A relation that metabolised \emph{every} misalignment, that returned instantly and completely to concordance after each passing, would be a relation incapable of the dormancy and return through which generative systems renew themselves---a relation, as the cosmological sections will argue (\xref{sec:dissipation}--\xref{sec:rose}), at equilibrium, and so dead. Some shadow is the price a living bond pays for remaining alive. The clinical art---and the practical art the closing sections reconstruct (\xref{sec:praxis})---begins in the capacity to tell which shadow is which.

The difficulty is that the two are not distinguishable \emph{by their surface}. A drifting-out-of-phase that is the ordinary respiration of a long bond, and a drifting that is the leading edge of domination, can present identically to the one living through them. This is precisely why a symptomatology is insufficient and an etiology is required: the difference between reparable and structural lies not in how the symptom looks but in \emph{how it is caused}---in which of the layers of the bond the disturbance originates, and whether its source is internal to the dyad or pressed upon it from the field of power without. To that question of causation the paper now turns.

\subsection{A note on what the symptoms share}
\label{sub:symptom-shared}

One observation closes the description and frames the diagnosis. The three presentations---estrangement, domination, the decoherence of trust---look like three different troubles, but each is, at bottom, a disturbance in the same thing: the capacity of two opaque subjects to go on \emph{generating} value between them rather than extracting it, settling it, or freezing it. Estrangement is generation arrested mid-cycle; domination is generation captured and made to flow one way; decoherence is the symbolic form of generation outliving its real substance. The common object of all three is the generative coupling itself; the common effect is its conversion, by one route or another, from a living circulation into something static---passed, captured, or hollow. The chapters that follow ask, layer by layer, how that conversion comes about, beginning with the most intimate and least culpable of its causes: the simple fact that two subjects do not speak, in the matter of trust, the same language.
